\section{Deductive explanation: complete proofs}
\label{sec:annex-p16-deductive-nomological-explanation}

Use the equality-free relational language and classical
consequence of \cref{def:p16-dn-language}.
The following proofs supply the logical results of the
1948 presentation, together with the explicitly stated
handbook examples.

\begin{annexitem}{annex:p16.dn-ground-language}{thm:p16-dn-ground-language}{Realizing ground assignments and separating general sentences}
\begin{proof}
Take as domain the set of constant names, with each
name denoting itself. If there are no names, add
one object. For each predicate and tuple of named
objects, assign membership according to the prescribed
truth value of the corresponding ground atom.
Distinct names give distinct tuples, so this prescription
is consistent. Assign arbitrary, for example false,
values on any remaining tuples. A nullary predicate,
if present, receives its prescribed truth value directly.
This realizes every ground assignment.

Truth of a singular sentence is determined recursively
by the Boolean truth clauses on those atoms.
Every first-order structure gives such an assignment,
and the preceding construction realizes every one.
Thus singular entailment agrees in both directions
with Boolean entailment.

To show essential generality, it suffices to exhibit
two structures agreeing on all ground atoms but
disagreeing on the proposed generalized sentence.
Interpret every constant as \(0\). For
\(\forall xP(x)\), use the domain \(\{0\}\)
with \(P(0)\) true, and then the domain
\(\{0,1\}\) with only \(P(0)\) true.
All ground atoms still concern named objects
denoting \(0\), while the universal changes truth.
Other predicates can agree on every named tuple.

For \(\forall x(P(x)\to Q(x))\), make both
predicates false at \(0\) in the one-object
structure. Add the unnamed object \(1\),
with \(P(1)\) true and \(Q(1)\) false.
The same ground assignment now extends to
a structure falsifying the universal.
For \(\forall xR(a,x)\), make \(R(0,0)\)
true in the one-object structure and keep
it true in the larger structure while
making \(R(0,1)\) false.
For \(\exists xP(x)\), make \(P(0)\) false,
then add the unnamed object \(1\) with
\(P(1)\) true. These pairs likewise agree
on every ground atom and disagree on
the stated sentence. No singular sentence
can therefore be equivalent to any of them.

If \(P(a)\) is true, every instance of
\(P(x)\lor P(a)\) is true. Conversely its
universal instance at the denotation of
\(a\) gives \(P(a)\lor P(a)\), hence
\(P(a)\). The other displayed universal
is true under every assignment at every
object, so is equivalent to \(\top\).
\end{proof}
\end{annexitem}

\begin{annexitem}{annex:p16.dn-law-hierarchy}{thm:p16-dn-law-hierarchy}{Every law is a true theory}
\begin{proof}
A fundamental law is purely universal and true,
hence purely generalized and true, so is a
fundamental theory. A derivative law is
universal and essentially generalized.
Because it is not purely universal, its
universal form contains an individual constant
and is not purely generalized. It follows
from fundamental laws, which are fundamental
theories. It is therefore a derivative theory.

A fundamental theory is true by definition.
For a derivative theory, the intended structure
satisfies each of the fundamental theories
from which it follows, and hence satisfies
the consequence. Thus every theory is true.
A sentence cannot at once be purely
generalized and not purely generalized;
the same observation applies to the
universal subclasses. This proves disjointness.
Purely universal logical truths satisfy
the definition. A universal conditional
with an empty antecedent extension is
also true in the intended structure,
whether or not its antecedent is
logically impossible.

Universal instantiation of
\(\forall y\forall xR(y,x)\) at \(y=a\)
gives \(\forall xR(a,x)\). The latter
is essentially universal by the preceding
lemma and contains \(a\), establishing
the derivative-law example.
\(\exists xP(x)\) has a purely generalized
existential prefix; when true it is a
fundamental theory, but its prefix is
not universal.

For the cancellation example choose the
two structures used above to vary
\(U=\forall xQ(x)\), with all ground
\(Q\)-atoms true in both, and make
\(P(a)\) false in both. Then \(A\lor U\)
changes truth with \(U\), and \(A\lor\neg U\)
changes in the opposite direction.
They are both essentially generalized.
Boolean distributivity gives
\[
 (A\lor U)\land(A\lor\neg U)
          \equiv A\lor(U\land\neg U)\equiv A .
\]
The conjunction is equivalent to a singular
sentence and is therefore not essentially
generalized, despite the essential generality
of its two constituents.
\end{proof}
\end{annexitem}

\begin{annexitem}{annex:p16.dn-verification-properties}{thm:p16-dn-verification-properties}{Consistency, gain and the two witness models}
\begin{proof}
Let \(\Gamma\) be a witness. A structure satisfying
\(\{T\}\cup\Gamma\) also satisfies \(C\), because
\(\Gamma\entails C\). It therefore witnesses
satisfiability of \(T\land C\).
If \(C\entails E\), transitivity with
\(\Gamma\entails C\) would give
\(\Gamma\entails E\), contrary to the witness
condition. Thus \(C\not\entails E\).

A valid \(E\) follows from every set of
premises, contradicting the same condition.
A contradictory \(E\), if entailed by
the satisfiable \(T\land C\), would be
true in its model, which is impossible.
Hence \(E\) is neither valid nor contradictory.
The failure of \(C\) alone to entail \(E\)
establishes the collective necessity of
the theoretical component in the stated
deduction; it says nothing about minimality
of every individual premise.

For a truth-qualified explanans, \(T\)
is true in \(\mathfrak A\) because it
is a theory, and \(C\) is true there
by definition. Soundness of the deduction
therefore gives \(\mathfrak A\models E\).

Logical equivalence preserves consequence,
truth in every structure and satisfiability
with any set of premises. Thus the same
\(\Gamma\) works after any of the stated
equivalent replacements, provided the
required singular syntactic forms are
retained. Essential generality is itself
invariant under equivalence.

Compatibility gives one model of
\(\{T\}\cup\Gamma\); failure of
\(\Gamma\entails E\) gives a model
of \(\Gamma\cup\{\neg E\}\).
These cannot be the same model:
every model of \(\{T\}\cup\Gamma\)
satisfies \(C\) and then \(E\).
This proves the distinction between
the two tests and the impossibility
of the strengthened joint countermodel
condition.
\end{proof}
\end{annexitem}

\begin{annexitem}{annex:p16.dn-finite-witness}{thm:p16-dn-finite-witness}{Reducing a witness to a finite literal cube}
\begin{proof}
Start with any witness \(\Gamma\).
It is consistent because \(\{T\}\cup\Gamma\)
has a model. Let \(\Gamma_0\) be those
literals of \(\Gamma\) whose atoms
occur in \(C\). There are at most
as many such literals as atoms of
\(C\), since consistency excludes
opposite literals on the same atom.

Suppose \(\Gamma_0\not\entails C\).
By the ground freedom lemma, there
is a Boolean assignment to the
atoms of \(C\) satisfying all of
\(\Gamma_0\) and falsifying \(C\).
Extend it to every other atom
mentioned in \(\Gamma\), giving
that atom its required literal
value. No conflict arises, since
\(\Gamma\) is consistent and all
overlap with the atoms of \(C\)
has already been included in
\(\Gamma_0\). Give remaining atoms
arbitrary values. The ground
realization lemma provides a
structure satisfying \(\Gamma\)
and falsifying \(C\), contradicting
\(\Gamma\entails C\). Hence
\(\Gamma_0\entails C\).

Any model of \(\{T\}\cup\Gamma\)
is a model of \(\{T\}\cup\Gamma_0\).
Also \(\Gamma_0\) cannot entail
\(E\), for its superset
\(\Gamma\) would then entail \(E\).
Thus \(\Gamma_0\) is the required
finite witness.

For \(m\) atoms there are \(3^m\)
ways to assert, deny or omit each
one, including the empty set.
Truth-table evaluation over the
finite union of the atoms of
\(C,E\) decides the two singular
consequence tests. The remaining
test contains \(T\), which can
be quantified. Enumerating finitely
many such tests does not supply
a decision procedure for an
arbitrary individual first-order
satisfiability query.
\end{proof}
\end{annexitem}

\begin{annexitem}{annex:p16.dn-common-content}{thm:p16-dn-common-content}{The explicit proof of source Theorem 7.7}
\begin{proof}
Boolean distributivity gives
\[
 (C\lor E)\land(\neg C\lor E)\equiv E,
             \qquad C\land(C\lor E)\equiv C .
\]
The premise \(T,C\entails E\)
is equivalent to \(T\entails C\to E\),
that is, \(T\entails\neg C\lor E\).
These are exactly the three required
relations for \(E_1,E_2,C_1\).
All three formulas are singular
because \(C,E\) are singular.

Both \(C\) and \(E\) entail
\(C\lor E\). If they both entail
some \(H\), every structure satisfying
their disjunction satisfies at
least one of them and therefore
satisfies \(H\). Thus their
disjunction entails every common
consequence and is the strongest
one in the logical consequence
order.

If there is no nonvalid common
consequence, \(C\lor E\) is valid.
Together with
\(T\entails\neg C\lor E\),
the first displayed identity now
gives \(T\entails E\).
This proves the consequence about
dispensability of the particular
conditions without imposing
any interpretation of actual
reasoning order.
\end{proof}
\end{annexitem}

\begin{annexitem}{annex:p16.dn-smuggled-result}{thm:p16-dn-smuggled-result}{Rejecting the source's conditional disguise}
\begin{proof}
Universal instantiation of \(T\)
gives \(P(a)\), and modus ponens
with \(C\) gives \(R(a,b)\).
A structure with \(P(a)\) false
and \(R(a,b)\) false satisfies
\(C\) and falsifies \(E\), so
\(C\not\entails E\). Essential
generality of \(T\) was proved
in the ground-language lemma.
A structure with all \(P\)'s
true and \(R(a,b)\) true makes
both proposed premises true.

If there were a final verification
witness, the finite-witness theorem
would give one using only the
atoms \(P(a),R(a,b)\).
A consistent literal set entails
\(\neg P(a)\lor R(a,b)\) only
if it includes \(\neg P(a)\)
or \(R(a,b)\). Indeed, when
neither is included, setting
\(P(a)\) true and \(R(a,b)\)
false is compatible with every
remaining literal and falsifies
that disjunction.
Compatibility with \(T\) rules
out \(\neg P(a)\), so the witness
must contain \(R(a,b)=E\).
It then entails \(E\), a
contradiction. This proves failure
of the final criterion.

For the general construction,
\(T\entails A\) and
\(C=A\to E\) give \(T,C\entails E\).
Any structure satisfying
\(\neg A\land\neg E\) satisfies
the conditional \(C\) while
falsifying \(E\), so \(C\not\entails E\).
When \(E\) is true in the intended
structure, the same conditional
is true there. If \(A\) is valid,
however, \(A\to E\) is equivalent
to \(E\), and the required
countermodel does not exist.
This verifies the exact
qualification rather than an
unrestricted claim about all
lawlike sentences.
\end{proof}
\end{annexitem}

\begin{annexitem}{annex:p16.dn-worked-certificates}{thm:p16-dn-worked-certificates}{Every compatibility and non-entailment check}
\begin{proof}
For the first case, take
\(\Gamma=\{P(a)\}\).
It entails \(C\).
A one-object structure with
both \(P,Q\) true satisfies
\(\{T\}\cup\Gamma\).
A structure with \(P(a)\)
true and \(Q(a)\) false
satisfies \(\Gamma\) but
not \(E\). Instantiation
and modus ponens give
\(T,C\entails E\).
The ground-language lemma
supplies essential generality.

For the second, the empty
\(\Gamma\) entails \(\top\),
is compatible with the
satisfiable \(\forall xP(x)\),
and does not entail \(P(a)\):
the all-false \(P\) structure
is a countermodel.
The universal directly entails
the result. This is an example
where the initial component
is wholly dispensable.

For the third, let
\(\Gamma=\{R(a,b),P(a),U(b)\}\).
It entails \(C\) and is
compatible with \(T\), for
example by making \(P,Q,U\)
and \(R\) universal in a
nonempty structure.
The same literals can instead
hold with \(Q(a)\) false,
showing that \(\Gamma\)
does not entail \(E\).
Using \(T\) derives \(Q(a)\);
the other conjunct \(R(a,b)\)
already belongs to \(C\).
Thus the final criterion
allows this shared content.

For the fourth, \(\{Q(a)\}\)
entails \(P(a)\lor Q(a)\)
and does not entail \(P(a)\),
as witnessed by \(Q(a)\)
true and \(P(a)\) false.
It is compatible with
\(\forall xP(x)\) in a
structure with both predicates
universal. The theoretical
component already entails
\(E=P(a)\).
Now choose the intended
structure with \(P\) universal
and \(Q\) empty. It makes
\(T,C,E\) true while the
chosen witness literal is
false. The witness requirement
is existential over compatible
verifications; it does not
require that witness to be
true in this particular
intended structure.
\end{proof}
\end{annexitem}

\begin{annexitem}{annex:p16.dn-emergence-monotonicity}{thm:p16-dn-emergence-monotonicity}{Theory and descriptive enlargement}
\begin{proof}
If \(\Gamma\subseteq\Gamma'\),
every model of \(\Gamma'\)
is a model of \(\Gamma\).
Therefore \(\Gamma\entails W(w)\)
implies \(\Gamma'\entails W(w)\).
Apply this with
\(\Gamma=\mathcal T\cup D_{\mathsf{Pt},G}(w)\)
and with any enlargement of
its theory or descriptive
premises.
Failure of consequence can
disappear under enlargement,
but a previously valid
consequence cannot disappear.

If an enlargement is
inconsistent, it has no
countermodels and entails
every formula. It cannot
be a collection of true
premises in the intended
structure, so this vacuous
consequence is not a sound
factual explanation.

Finally suppose \(p\) is
a named part and the
available data include
\(\mathsf{Pt}(p,w)\land W(w)\)
as a part attribute.
Conjunction elimination gives
\(W(w)\) without any theory.
Thus allowing this
whole-dependent attribute
eliminates non-derivability.
With no part there is no
such supplied instance,
so the argument cannot be
extended to a partless case
merely by an empty universal
description.
\end{proof}
\end{annexitem}

\begin{annexitem}{annex:p16.dn-emergence-examples}{thm:p16-dn-emergence-examples}{Three explicit part-to-whole models}
\begin{proof}
Use the three-element domain
\(\{p,q,w\}\) with each
name denoting itself,
\[
 \mathsf{Pt}=\{(p,w),(q,w)\},
                       \qquad A=\{p,q\}.
\]
First set \(R=\varnothing\)
and \(W=\varnothing\).
The displayed bridge theory
is true because its antecedent
is never satisfied, and the
initial part and attribute
data hold. Yet \(W(w)\)
is false. This is a
countermodel to the initial
deduction.

An intended occurrence can
instead be represented by
the same part and attribute
interpretations with
\(R=\{(p,q)\}\) and
\(W=\{w\}\). The bridge
holds, and \(W(w)\) is
true. Initially omitting
the connection fact from
the available description
leaves the preceding
countermodel possible.
Adding \(R(p,q)\) to
the data supplies every
antecedent conjunct of
the bridge instance
at \(x=p,y=q,z=w\);
modus ponens yields \(W(w)\).

If the bridge theory is
removed, keep all enlarged
data, including
\(R=\{(p,q)\}\), but
set \(W=\varnothing\).
Those data are still true,
while \(W(w)\) is false.
Thus knowing the connection
and knowing a rule that
connects it to the whole
are distinct requirements.
\end{proof}
\end{annexitem}

\begin{annexitem}{annex:p16.dn-thermal-model}{thm:p16-dn-thermal-model}{The complete drop-and-rise calculation}
\begin{proof}
Multiplying \(\dot u+au=a\)
by \(e^{at}\) gives
\((e^{at}u)'=ae^{at}\).
Integration from zero to
\(t\), with the initial
condition, yields
\(u=1-e^{-at}\).
The same calculation
gives \(v=1-e^{-bt}\).
For any two solutions
their difference satisfies
\(\dot z+az=0,z(0)=0\);
the integrating factor
makes \(e^{at}z\)
constant zero. This
proves uniqueness, and
the analogous argument
applies to \(v\).

Substitution into the
composition equation
gives the stated \(H\).
Its derivative is
\[
 H'(t)=e^{-bt}
           \bigl(Aa\,e^{(b-a)t}-Bb\bigr).
\]
The factor outside
parentheses is positive.
The expression inside
is strictly increasing,
starts negative by
\(Bb>Aa\), and has
its unique zero at
\(t_*\) as stated.
Thus \(H\) strictly
decreases from \(H(0)=0\)
to its negative minimum
and then strictly
increases. Its limit
is \(A-B>0\).
Continuity and strict
increase after the
minimum give exactly
one later zero.

For the numerical
parameters, set \(z=e^{-t}\).
The equation for a zero
is
\[
 1-2z+z^3=(z-1)(z^2+z-1)=0 .
\]
Besides \(z=1\), corresponding
to \(t=0\), the sole root
in \((0,1)\) is
\((\sqrt5-1)/2\).
Taking the negative
logarithm gives
\(t_0=\log((1+\sqrt5)/2)\).
The assumptions are
substantive: for
\(A=2,B=1,a=b=1\),
one obtains
\(H(t)=1-e^{-t}\),
which increases from
the start and has
no initial drop.
\end{proof}
\end{annexitem}

\begin{annexitem}{annex:p16.dn-function-and-motive}{thm:p16-dn-function-and-motive}{The missing logical directions}
\begin{proof}
The valuation \(B=V=\mathbf F\)
satisfies \(B\to V\)
while falsifying both
of the proposed categorical
conclusions. For the
alternative-realization case,
take \(V=B'=\mathbf T\)
and \(B=\mathbf F\).
Then \(V\) and
\(V\to(B\lor B')\)
are true while \(B\)
is false.

For the motivational
case, the three current
conditions give
\(D\land K\land O\).
The stipulated conditional
therefore yields \(A\)
by modus ponens.
Set \(D,K,O,A\)
true and \(G\)
false. This valuation
satisfies all the
premises but not the
future achievement.
Thus the deduction
of current action
does not require
truth of that future
event. Whether the
illustrative action
rule is empirically
adequate is an
additional question.
\end{proof}
\end{annexitem}

\begin{annexitem}{annex:p16.dn-intervention-boundary}{thm:p16-dn-intervention-boundary}{The two observationally identical causal models}
\begin{proof}
Use \(\Omega=\{0,1\}\)
with probability one
half at each point,
and \(U(\omega)=\omega\).
The first model has
the acyclic edge
\(X\to Y\); the
second \(Y\to X\).
Their equations give
\((X,Y)=(U,U)\)
in every ordinary
context. Hence their
observational laws
are identical, each
assigning one half
to \((0,0)\) and
one half to \((1,1)\).

After replacing only
the \(X\) equation
by the constant zero,
the first model
still uses \(Y=X\),
so gives \((0,0)\)
in every context.
The second retains
\(Y=U\), giving
\((0,0)\) and
\((0,1)\) with equal
probabilities.
Thus the intervention
laws differ even
though the ordinary
constraint and every
observational probability
agree. Deduction
within the common
observational relation
does not determine
which equation is
retained under the
intervention.
\end{proof}
\end{annexitem}
